\scanpage{156}
% archive foliation: 155
qu'il est \uline{surjectif} et \uline{radiciel}, donc que pour tout $y\in Y$, la fibre géométrique de $P/H$ \deleted{\unsure{dans}} en $y$ s'envoie bijectivement sur celle de $X$ en $y$.
Or c'est encore un pt trivial, d'après l'explication précédente.

\margin{pts géométriques \unsure{en $y$} de $P/H$ $=$ \deleted{\unsure{pts géom.}} \unsure{les classes} \unsure{sous $H$} des pts géom. de $P$.}

Évident $(iii\ \mathrm{bis})\Rightarrow(iii)$, et $(iii)\Rightarrow(i)$ d'après un \uline{lemme} précédent, déjà utilisé.
D'où l'équivalence de $(i)\ (iii\ \mathrm{bis})\ (iii)$.

De plus $(iii)\deleted{\unsure{\Leftarrow}}\Rightarrow(ii\ \mathrm{bis})$, comme on voit en prenant $Y_1=P$, et utilisant la formule
$(X\times G)/H\simeq X\times G/H$, et $(ii\ \mathrm{bis})\Rightarrow(ii)$.
Enfin, $(ii\deleted{\mathrm{bis}})\Rightarrow(i)$ est immédiat
\deleted{[car la condition \unsure{loc.}}
Car si $X$ est défini par un faisceau d'alg $\underline A$ sur $Y$, $X_1$ le sera par son image réciproque $\underline A_1$.
Si $\underline A_1$ est loc. libre, il en est de $=$ de $\underline A$ comme bien connu.
Si les fibres de $X_1$ sur $Y_1$ sont séparables, \inserted{(déduites de celles de $X$ sur $Y$ par ext. des corps de base)} il en est de $=$ de celles de $X$ sur $Y$.
\uline{Donc} si $X_1$ est étale sur $Y_1$, $X$ l'est sur $Y$.
\margin{\uline{loc.}}

\uline{Cor}
Si $X$ est connexe, alors dans $(iii)$ on peut prendre $P$ \uline{connexe}.
\figure{\begin{tikzpicture}[x=1cm,y=1cm]
\draw (0,.65) -- (0,-.3) .. controls (0,-.45) and (.04,-.5) .. (.08,-.55);
\end{tikzpicture}}
